\apphead{Appendix 1}{Motivations}

In this "motivational" appendix we will first recall the relation between
complex elliptic curves\index{elliptic curve} and lattices in $\C$, then the relation between modular
forms and the de Rahm cohomology of elliptic curves, and finally the relation
between the Gauss-Manin connection\index{Gauss--Manin connection} and Serre's $\partial$ operator on modular
forms. These relations are due to Weierstrass and Deligne.

\sechead{A1.1}{Lattices and elliptic curves}

Given a lattice $L\subset\C$, we may form the quotient $\C/L$, a one-dimensional
complex torus, and endow it with the translation-invariant one-form $\omega=\dd z$
($z$ the coordinate on $\C$). Thanks to Weierstrass, we know that $\C/L$
\uline{is} an elliptic curve\index{elliptic curve}, given as a cubic $\PP^2$ by the inhomogeneous
equation

\begin{kadispb}{A1.1.1}
y^2 = 4x^3 - g_2x - g_3 ,
\end{kadispb}

such that $\omega$ is the differential $\dd x/y$. The isomorphism from $\C/L$
to this curve is explicitly given by the $\wp$-function:

\begin{kadispb}{A1.1.2}
z\in\C/L \longrightarrow (x=\wp(z;L),\ y=\wp'(z;L))
\end{kadispb}

where

\begin{kadispb}{A1.1.2.1}
\wp(z;L) = \frac{1}{z^2} + \sum_{\ell\in L-\{0\}}
\left(\frac{1}{(z-\ell)^2} - \frac{1}{\ell^2}\right) ,
\end{kadispb}

\begin{kadispb}{A1.1.2.2}
\wp'(z;L) = \frac{\dd\,\wp(z;L)}{\dd z} = \frac{-2}{z^3} +
\sum_{\ell\in L-\{0\}} \frac{-2}{(z-\ell)^3} ,
\end{kadispb}

\begin{kadispb}{A1.1.2.3}
g_2 = 60\sum_{\ell\in L-\{0\}} 1/\ell^4 ,\quad
g_3 = 140\sum_{\ell\in L-\{0\}} 1/\ell^6 .
\end{kadispb}

Conversely, given an elliptic curve $E$ over $\C$ together with a non-zero
everywhere holomorphic differential $\omega$, it arises in the above way from the
lattice of periods of $\omega$,

\begin{kadispb}{A1.1.2.4}
L(E,\omega) = \left\{\int_\gamma \omega \mid \gamma\in H_1(E;\Z)\right\}\subset\C .
\end{kadispb}

Under this correspondence, the effect of replacing $(E,\omega)$ by
$(E,\lambda\omega)$, $\lambda\in\C^*$, is to replace $L$ by $\lambda\cdot L$:

\begin{kadispb}{A1.1.2.5}
L(E,\lambda\omega) = \lambda\cdot L(E,\omega) .
\end{kadispb}

Recall that classically, a complex modular form\index{modular form} of weight $k$ (and level 1) is a
holomorphic function on the upper-half plane $f(\tau)$ which satisfies the
transformation equation

\begin{kadispb}{A1.1.3}
f\left(\frac{a\tau+b}{c\tau+d}\right) = f(\tau)\cdot(c\tau+d)^k
\qquad\forall\left(\begin{smallmatrix}a&b\\ c&d\end{smallmatrix}\right)\in SL_2(\Z) .
\end{kadispb}

As explained in \cite{serre70}, there is associated to $f$ a unique function
of \uline{lattices} $F(L)$ such that $f(\tau)=F(\Z\tau+\Z)$, and which is
homogeneous of degree $-k$ in $L$: $F(\lambda L)=\lambda^{-k}F(L)$ for
$\lambda\in\C^*$. (Explicitly, $F(L)=\omega_2^{-k}f(\omega_1/\omega_2)$ if
$L=\Z\omega_1+\Z\omega_2$ and $\operatorname{Im}(\omega_1/\omega_2)>0$.)

By Weierstrass, we may now associate to $f$ a "holomorphic" function
$\mathbb{F}$ of pairs $(E,\omega)$ consisting of an elliptic curve$/\C$ together
with a nowhere-vanishing differential which is homogeneous of degree $-k$ in the
second variable: $\mathbb{F}(E,\lambda\omega)=\lambda^{-k}\mathbb{F}(E,\omega)$,
defined by $\mathbb{F}(E,\omega)=F(L(E,\omega))$. This is the point of view
taken in the text.

\sechead[A1.2a]{A1.2}{Homomorphy at $\infty$ and the Tate curve}

Recall further that a complex modular form\index{modular form} $f(\tau)$ is said to be
\uline{meromorphic} (resp. holomorphic) at $\infty$, if the periodic function
$f(\tau)=f(\tau+1)$, when viewed as a function of $q=\exp(2\pi i\tau)$,
holomorphic for $0<|q|<1$, in fact extends to a meromorphic (resp. holomorphic)
function of $q$ in $|q|<1$. In terms of $\mathbb{F}$, we are asking about the
behavior of

\[
\mathbb{F}\left(\C/2\pi i\Z+2\pi i\tau\Z,2\pi i\,\dd z\right) =
\mathbb{F}\left(\C^*/q^\Z,\dd t/t\right)
\]

(where $t=\exp(2\pi iz)$ is the parameter on $\C^*$, and $q^\Z$ denotes the
subgroup of $\C^*$ generated by $q$), as $q$ tends to zero. By standard
calculations (cf. \cite{roquette70}), the curve $\C/L$, $L=2\pi i\Z+2\pi i\tau\Z$ with
differential $2\pi i\,\dd z$ is given as the plane cubic

\[
Y^2 = 4X^3 - \frac{E_4}{12}X + \frac{E_6}{216} ,\quad\text{with differential}\ \dd X/Y
\]

\begin{kadispb}{A1.2.1}
(X = \wp(2\pi iz,L),\quad Y = \wp'(2\pi iz,L))
\end{kadispb}

with coefficients the Eisenstein series

\begin{kadispb}{A1.2.2}
\begin{cases}
12\cdot(2\pi i)^4 g_2(\tau) = E_4 = 1 + 240\sum\sigma_3(n)q^n \\
216\cdot(2\pi i)^6 g_3(\tau) = E_6 = 1 - 504\sum\sigma_5(n)q^n
\end{cases}
;\quad \sigma_k(n) = \sum_{\substack{d\mid n\\ d\geq 1}} d^k
\end{kadispb}

Thus to ask that the modular form $f$ be meromorphic (resp. holomorphic) at
$\infty$ is to ask that
$\mathbb{F}\left(Y^2=4X^3-\frac{E_4}{12}X+\frac{E_6}{216},\ \dd X/Y\right)$ lie
in the ring $\C((q))$ of finite-tailed Laurent series (resp., that it lie in
$\C[[q]]$, the ring of formal power series in $q$).

The equation \ref{ka:A1.2.1} in fact defines an elliptic curve\index{elliptic curve} over the ring
$\Z[1/6]((q))$; in fact, if we introduce

\[
X = x + \frac{1}{12} ,\qquad Y = x + 2y
\]

then we may rewrite the equation in the form

\begin{kadispb}{A1.2.3}
y^2 + xy = x^3 + \mathbb{B}(q)x + \mathbb{C}(q)
\end{kadispb}

with coefficients

\begin{kadispb}{A1.2.4}
\begin{cases}
\mathbb{B}(q) = -5\left(\frac{E_4-1}{240}\right) = -5\sum_{n\geq 1}\sigma_3(n)q^n \\
\mathbb{C}(q) = \frac{-5\left(\frac{E_4-1}{240}\right)-7\left(\frac{E_6-1}{-504}\right)}{12}
= \sum_{n\geq 1}\left(\frac{-5\sigma_3(n)-7\sigma_5(n)}{12}\right)q^n .
\end{cases}
\end{kadispb}

\uline{This last equation} defines an elliptic curve over $\Z((q))$ whose
restriction to $\Z[\frac{1}{6}]((q))$ is the above curve, and the nowhere
vanishing differential $\dd x/2y+x$ restricts to give $\dd X/Y$ over
$\Z[\frac{1}{6}]((q))$.

By definition, the Tate curve\index{Tate curve} $\mathrm{Tate}(q)$ with its \uline{canonical
differential} $\omega_{\mathrm{can}}$ is the elliptic curve over $\Z((q))$
defined by \eqref{ka:A1.2.3}, with differential $\omega_{\mathrm{can}}=\dd x/2y+x$. For
each integer $n\geq 1$, the Tate curve $\mathrm{Tate}(q^n)$ with its canonical
differential $\omega_{\mathrm{can}}$ is deduced from
$(\mathrm{Tate}(q),\omega_{\mathrm{can}})$ by the extension of scalars
$\Z((q))\to\Z((q))$ given by $\Sigma\,a_iq^i\to\Sigma\,a_iq^{ni}$. Explicitly,
$(\mathrm{Tate}(q^n),\omega_{\mathrm{can}})$ is given by

\begin{kadispb}{A1.2.5}
y^2+xy = x^3+\mathbb{B}(q^n)\cdot x+\mathbb{C}(q^n) ;\quad
\omega_{\mathrm{can}} = \dd x/2y+x .
\end{kadispb}

Let $\zeta_n$ be a primitive $n$'th root of unity. The points of order $n$ on
$\C^*/q^{n\Z}$ are clearly the (images of the) $n^2$ points

\begin{kadispb}{A1.2.6}
(\zeta_n)^iq^j ,\qquad 0\leq i,j\leq n-1 .
\end{kadispb}

Using the \uline{explicit expressions} for $x$ and $y$ as functions of
$t=\exp(2\pi iz)$

\begin{kadispb}{A1.2.7}
\begin{cases}
x(t) = \sum_{k\in\Z}\frac{q^{nk}t}{(1-q^{nk}t)^2} - 2\sum_{k=1}^{\infty}\frac{q^{nk}}{1-q^{nk}} \\
y(t) = \sum_{k\in\Z}\frac{(q^{nk}t)^2}{(1-q^{nk}t)^3} + \sum_{k=1}^{\infty}\frac{q^{nk}}{1-q^{nk}} ,
\end{cases}
\end{kadispb}

one sees that each of the $n^2-1$ points $(\zeta_n)^iq^j$,
$0\leq i,j\leq n-1$, $(i,j)\neq(0,0)$ has $x$ and $y$ coordinates in
$\Z[[q]]\tens_\Z\Z[\zeta_n,1/n]$. Hence \uline{all} level $n$ structures on
$\mathrm{Tate}(q^n)$ over $\Z((q))$ are defined over
$\Z((q))\tens\Z[\zeta_n,1/n]$ (rather than just over
$\Z[\zeta_n,1/n]((q))$). This implies that the q-expansions of a modular form
of level $n$ have \uline{bounded denominators} (cf.\eqref{ka:1.2.1}).

\sechead{A1.2}{Modular forms and de Rham cohomology}

We can now give a purely algebraic definition of modular forms\index{modular form} of weight $k$,
(meromorphic at $\infty$) as being certain "functions" $f(E,\omega)$ defined
whenever

\[
\begin{tikzcd}
E\ ;\ \omega \arrow[d, "\pi"'] \\
R
\end{tikzcd}
\]

is any elliptic curve\index{elliptic curve} over any ring $R$, and $\omega\in\Gamma(E,\Omega^1_{E/R})$
is a nowhere vanishing differential on $E$, whose \uline{values} $f(E,\omega)$
are elements of the ground-ring $R$. The conditions to be satisfied are

\kaitemb{0)} $f(E,\lambda\omega)=\lambda^{-k}f(E,\omega)$ for all
$\lambda\in R^\times$ ;

\kaitemb{1)} $f(E,\omega)$ depends only on the \uline{isomorphism class} of
$(E,\omega)$ over $R$ ;

\kaitemb{2)} if $\varphi\colon R\to R'$ is a ring homomorphism, then, denoting by
$(E_\varphi,\omega_\varphi)$ the curve with differential over $R'$ deduced by
extension of scalars, we have $f(E_\varphi,\omega_\varphi)=\varphi(f(E,\omega))$.

\{Such modular forms are \uline{automatically} meromorphic at infinity, simply
because the Tate curve\index{Tate curve} $\mathrm{Tate}(q)$ \uline{is} an elliptic curve over
$\Z((q))$.\}

Given a modular form $f$ of weight $k$, we may form the $k$-ple differential
$f(E,\omega)\cdot\omega^{\otimes k}$ on $E$, which is \uline{independent} of the
choice of $\omega$, and view it as a global section over $R$ of the (invertible)
sheaf $(\w_{E/R})^{\otimes k}$, where

\[
\w_{E/R} \xlongequal{\text{dfn}} \pi_*(\Omega^1_{E/R}) .
\]

This permits us to interpret a (meromorphic at $\infty$) modular form of weight
$k$ as a function $f(E)$, defined on any elliptic curve $E$ over any ring $R$,
with values in the global sections of $(\w_{E/R})^{\otimes k}$, which satisfies

\kaitemb[1-b]{1)} if $\alpha\colon E\to E'$ is an isomorphism of elliptic curves over
$R$, then $\alpha^*(f(E'))=f(E)$ ;

\kaitemb[2-b]{2)} if $\varphi\colon R\to R'$ is a ring homomorphism, then
$f(E_\varphi)=\varphi^*(f(E))$.

Why bother to look at the de Rham cohomology\index{de Rham cohomology}? Over \uline{any} base ring $R$,
the (1\uline{st}) de Rham cohomology of an elliptic curve $E/R$, noted
$H^1_{\dR}(E/R)$ and defined as $\mathbb{H}^1(E,\widehat{\Omega}^\bullet_{E/R})$,
sits in a short exact sequence, its "Hodge filtration") of $R$-modules

\begin{kadispb}[A1.2.1-b]{A1.2.1}
0 \longrightarrow \w_{E/R} \longrightarrow H^1_{\dR}(E/R) \longrightarrow \Lie(E/R) \longrightarrow 0 \\
\qquad\qquad\qquad\qquad\Vert \\
\qquad\qquad\qquad H^1(E,\widehat{\cO}_E) \\
\qquad\qquad\qquad\qquad\Vert \\
\qquad\qquad\qquad (\w_{E/R})^{\otimes -1}
\end{kadispb}

Furthermore, when the integer $6$ is invertible in $R$, this sequence has a
\uline{canonical} (but not \uline{functorial}) splitting,

\[
H^1_{\dR}(E/R) \longrightarrow \w_{E/R}\oplus(\w_{E/R})^{\otimes-1}
\]

which may be obtained as follows. Given $(E,\omega)$ over $R$, $R\ni 1/6$, then
there are unique meromorphic functions with poles only along the identity
section, of orders 2 and 3 respectively, $X$ and $Y$ on $E$ such that
$\omega=\dd X/Y$ and such that $E$ is defined by (inhomogeneous) equation

\[
Y^2 = 4X^3 - g_2X - g_3 ,\qquad g_2, g_3\in R
\]

(when $R=C$, we have $X=\wp(z;L)$, $Y=\wp'(z;L)$, $L$ the lattice of periods of
$\omega$). To specify the dependence on $\omega$, let's write $X(E,\omega)$,
$Y(E,\omega)$, $g_2(E,\omega)$, $g_3(E,\omega)$. By uniqueness, we necessarily
have

\begin{kadispb}[A1.2.2-b]{A1.2.2}
\begin{cases}
X(E,\lambda\omega) = \lambda^{-2}\cdot Y(E,\omega) \\
Y(E,\lambda\omega) = \lambda^{-3}\cdot Y(E,\omega) \\
g_2(E,\lambda\omega) = \lambda^{-4}\cdot g_2(E,\omega) \\
g_3(E,\lambda\omega) = \lambda^{-6}\cdot g_3(E,\omega)
\end{cases}
\end{kadispb}

But over \uline{any} base-ring $R$, the first de Rham cohomology of an elliptic
curve $E/R$ is nothing other than the module of differentials on $E/R$ having
at worst double poles at $\infty$ (i.e. along the identity section). More
precisely, the \uline{inclusion} of the de Rham complex
$\widehat{\Omega}^\bullet_E\to\Omega^1_{E/R}$ in the complex

\[
\widehat{\cO}_E(\infty)\longrightarrow\Omega^1_{E/R}(2\infty)
\]

induces an \uline{isomorphism} on $H^1$. Because
$H^i(E,\widehat{\cO}_E(\infty))=0={}\allowbreak H^i(E,\Omega^1_{E/R}(2\infty))=0$ for $i>0$,
we have

\begin{kadispb}[A1.2.3-b]{A1.2.3}
\mathbb{H}^1(E,\widehat{\cO}_E(\infty)\to\Omega^1_{E/R}(2\infty))
&= \Coker(H^0(E,\widehat{\cO}_E(\infty))\to H^0(E,\Omega^1_{E/R}(2\infty))) \\
&= \Coker(R\xrightarrow{0}H^0(E,\Omega^1_{E/R}(2\infty))) \\
&= H^0(E,\Omega^1_{E/R}(2\infty)) .
\end{kadispb}

If we suppose $6$ to be \uline{invertible} in $R$, then \uline{as soon as} we
choose a nowhere vanishing differential $\omega$ on $E$, we may canonically
specify a basis of $H^1(E,\Omega^1_{E/R}(2\infty))$, namely

\begin{kadispb}[A1.2.4-b]{A1.2.4}
\omega = \frac{\dd X(E,\omega)}{Y(E,\omega)}\quad\text{and}\quad
\eta = X(E,\omega)\cdot\omega = \frac{X(E,\omega)\cdot\dd X(E,\omega)}{Y(E,\omega)} .
\end{kadispb}

Replacing $\omega$ by $\lambda\omega$, $\lambda\in R^\times$, has the effect of
replacing this basis by

\begin{kadispb}[A1.2.5-b]{A1.2.5}
\lambda\omega = \frac{\dd X(E,\lambda\omega)}{Y(E,\lambda\omega)}\quad\text{and}\quad
\lambda^{-1}\eta = \frac{X(E,\lambda\omega)\,\dd X(E,\lambda\omega)}{Y(E,\lambda\omega)} ,
\end{kadispb}

which is to say that we have defined an isomorphism

\begin{kadispb}[A1.2.6-b]{A1.2.6}
H^1_{\dR}(E/R) \longleftarrow \w_{E/R}\oplus\w^{-1}_{E/R}
\end{kadispb}

given locally on $R$ in terms of the choice of a nowhere vanishing $\omega$ by

\[
a\omega+b\eta \longleftrightarrow a\omega\oplus b\omega^{-1} .
\]

For every integer $k\geq 1$, the $k$'th symmetric power of this isomorphism
provides an isomorphism

\begin{kadispb}[A1.2.7-b]{A1.2.7}
\Sym^k(H^1_{\dR}(E/R)) \simeq (\w_{E/R})^{\otimes k}\oplus
(\w_{E/R})^{\otimes k-2}\oplus\dots\oplus
(\w_{E/R})^{\otimes-k} .
\end{kadispb}

\sechead{A1.3}{The Gauss-Main connection, and the function P}

We begin by computing the Gauss-Manin connection\index{Gauss--Manin connection} on $H^1_{\dR}(E/R)$ in the case
where $R$ is the ring of holomorphic functions of $\tau$, and $E$ is the relative
elliptic curve\index{elliptic curve} defined by the lattice $\Z+\Z\tau$. The \uline{dual} $H_1(E/R)$ of
$H^1_{\dR}(E/R)$ is $R$-free on the two families of paths $\gamma_1$ and
$\gamma_2$:

\begin{equation}\tag{A1.3.1}\label{ka:A1.3.1}
\begin{tikzcd}
& \tau \\
0 \arrow[ur, "\gamma_1"] \arrow[r, "\gamma_2"'] & 1
\end{tikzcd}
\qquad\text{on } E_\tau = \C/\Z+\Z\tau
\end{equation}

The Gauss-Manin connection in this context is the action
$\nabla_\tau=\nabla\left(\frac{\dd}{\dd\tau}\right)$ of $\frac{\dd}{\dd\tau}$ on
$H^1_{\dR}(E/R)$ given by the formula (cf. \cite{katz72}, 4.1.2)

\begin{kadisp}{A1.3.2}
\int_{\gamma_i}\nabla_\tau(\xi) = \frac{\dd}{\dd\tau}\int_{\gamma_i}\xi
\qquad\text{for }\xi\in H^1_{\dR}(E/R),\text{ and }i=1,2
\end{kadisp}

(i.e., it is the dual of the connection on $H_1(E/R)$ for which $\gamma_1$ and
$\gamma_2$ are the horizontal sections).

To actually compute, let's note by $\omega$ (resp. $\eta$) the cohomology
classes of $\frac{\dd x}{y}$ and $\frac{x\,\dd x}{y}$ respectively, and denote by
$\omega_i$, $i=1,2$ and $\eta_i$, $i=1,2$ the periods $\int_{\gamma_i}\omega$ and
$\int_{\gamma_i}\eta$, which we view simply as elements of $R$. We will
\uline{also} denote by $\gamma_1$ and $\gamma_2$ the elements of $H^1_{\dR}(E/R)$
defined by Poincar\'e duality and the requirement that for any
$\xi\in H^1_{\dR}(E/R)$, $\int_{\gamma_i}\xi=\langle\xi,\gamma_i\rangle$. Thus
$\langle\gamma_2,\gamma_1\rangle=1=-\langle\gamma_1,\gamma_2\rangle$, and
$\langle\gamma_1,\gamma_1\rangle=\langle\gamma_2,\gamma_2\rangle=0$. We have
$\omega_i=\langle\omega,\gamma_i\rangle$ and
$\eta_i=\langle\eta,\gamma_i\rangle$ for $i=1,2$. Hence we necessarily have

\begin{kadisp}{A1.3.3}
\begin{cases}
\omega = \omega_1\gamma_2 - \omega_2\gamma_1 \\
\eta = \eta_1\gamma_2 - \eta_2\gamma_1
\end{cases}
;\qquad
\begin{pmatrix}\omega_1 & -\omega_2 \\ \eta_1 & -\eta_2\end{pmatrix}
\begin{pmatrix}\gamma_2 \\ \gamma_1\end{pmatrix} = \begin{pmatrix}\omega \\ \eta\end{pmatrix}
\end{kadisp}

(because both sides have the same periods over both $\gamma_1$ and $\gamma_2$).

But the classical "period relation" of Legendre

\begin{kadisp}{A1.3.4}
\eta_1\omega_2 - \eta_1\omega_2 = 2\pi i
\end{kadisp}

[which expresses that the topological cup-product $\langle\omega,\eta\rangle$ is
$2\pi i$, or equivalently that the \uline{DR}-cup-product
$\langle\omega,\eta\rangle_{\dR}=1$]. This allows us to express $\omega$ and
$\eta$ in terms of $\gamma_1$ and $\gamma_2$:

\begin{kadisp}{A1.3.5}
2\pi i\begin{pmatrix}\gamma_2 \\ \gamma_1\end{pmatrix} =
\begin{pmatrix}-\eta_2 & \omega_2 \\ -\eta_1 & \omega_1\end{pmatrix}
\begin{pmatrix}\omega \\ \eta\end{pmatrix} .
\end{kadisp}

Applying $\nabla_\tau$, we \uline{annihilate} $\gamma_1$ and $\gamma_2$, hence,
noting $\frac{\dd}{\dd\tau}$ by a prime $'$,

\begin{kadisp}{A1.3.6}
0 = \begin{pmatrix}-\eta_2' & \omega_2' \\ -\eta_1' & \omega_1'\end{pmatrix}
\begin{pmatrix}\omega \\ \eta\end{pmatrix} +
\begin{pmatrix}-\eta_2 & \omega_2 \\ -\eta_1 & \omega_1\end{pmatrix}
\begin{pmatrix}\nabla_\tau(\omega) \\ \nabla_\tau(\eta)\end{pmatrix}
\end{kadisp}

an equation we may solve using Legendre's relation:

\begin{kadisp}{A1.3.7}
\begin{pmatrix}\nabla_\tau(\omega) \\ \nabla_\tau(\eta)\end{pmatrix}
&= \frac{-1}{2\pi i}\begin{pmatrix}\omega_1 & -\omega_2 \\ \eta_1 & -\eta_2\end{pmatrix}
\begin{pmatrix}-\eta_2' & \omega_2' \\ -\eta_1' & \omega_1'\end{pmatrix}
\begin{pmatrix}\omega \\ \eta\end{pmatrix} \\
&= \frac{-1}{2\pi i}\begin{pmatrix}\eta_1'\omega_2-\eta_2'\omega_1 & \omega_1\omega_2'-\omega_2\omega_1' \\
\eta_2\eta_1'-\eta_1\eta_2' & \eta_1\omega_2'-\eta_2\omega_1'\end{pmatrix}
\begin{pmatrix}\omega \\ \eta\end{pmatrix} .
\end{kadisp}

At this point we must recall that $\omega_1=\tau$, $\omega_2=1$ and Legendre's
relation becomes: $\eta_1-\tau\eta_2=2\pi i$. Fed back into \eqref{ka:A1.3.7}, this
information gives

\begin{kadisp}{A1.3.8}
\begin{pmatrix}\nabla_\tau(\omega) \\ \nabla_\tau(\eta)\end{pmatrix}
= \frac{-1}{2\pi i}\begin{pmatrix}\eta_2 & -1 \\
(\eta_2)^2-2\pi i\eta_2' & -\eta_2\end{pmatrix}\begin{pmatrix}\omega \\ \eta\end{pmatrix} .
\end{kadisp}

\begin{kathmstar}{Lemma A1.3.9.}
$\eta_2 = -\Sigma_m\Sigma_n'\frac{1}{(m\tau+n)^2} = \frac{-\pi^2}{3}P$, where
$\Sigma'$ means that the term $(m=0,n=0)$ is omitted, and $P$ is the function of
$q=e^{2\pi i\tau}$ given by $P(q)=1-24\sum_{n\geq 1}\sigma_1(n)q^n$, where
$\sigma_1(n)=\sum_{d\geq 1,d\mid n}d$.
\end{kathmstar}

\begin{kaproof}
The first follows from the definition of $\eta_2$ as a \uline{period} of
$\eta=X\,\dd X/Y=\wp'(z)\dd z$, and the fact that $\eta=-\dd\zeta$, where $\zeta$
is the Weierstrass $\zeta$-function

\begin{kadisp}{A1.3.10}
\zeta(z) = \frac{1}{z} + \sum_m\sum_n'\left\{
\frac{1}{z-m\tau-n} + \frac{1}{m\tau+n} + \frac{z}{(m\tau+n)^2}\right\} .
\end{kadisp}

Indeed $\eta_2=\int_{\gamma_2}\eta = \int_0^1(-\dd\zeta(z)) = \int_z^{z+1}(-\dd\zeta(z))
= \zeta(z)-\zeta(z+1)$, and hence

\begin{kadisp}{A1.3.11}
\eta_2 &= \frac{1}{z}-\frac{1}{z+1} + \sum_m\sum_n'\left\{
\frac{1}{z-m\tau-n} - \frac{1}{z-m\tau-n+1} - \frac{1}{(m\tau+n)^2}\right\} \\
&= \frac{1}{z}-\frac{1}{z+1} + \sum_{m\neq 0}\sum_n\frac{-1}{(m\tau+n)^2}
+ \sum_{n\neq 0}\left\{\frac{-1}{n^2}+\frac{1}{z-n}-\frac{1}{z+1-n}\right\} \\
&= \sum_m\sum_n'\frac{1}{(m\tau+n)^2} .
\end{kadisp}

The second equality is ubiquitous (cf. \cite{serre70}, pp. 154-155).
\end{kaproof}

\begin{kathmstar}{Remark A1.3.12.}
A similar calculation, based on the fact that the $\zeta$ is an absolutely
convergent double sum, hence also given by function

\begin{kadisp}{A1.3.12.1}
\zeta(z) = \frac{1}{z} + \sum_n\sum_m'\left\{
\frac{1}{z-m\tau-n} + \frac{1}{m\tau+n} + \frac{z}{(m\tau+n)^2}\right\}
\end{kadisp}

shows that $\eta_1=\zeta(z)-\zeta(z+\tau)=-\Sigma_n\Sigma_m'\frac{\tau}{(m\tau+n)^2}$.
Comparing these two formulas, we see that $\eta_2(-1/\tau)=\tau\eta_1(\tau)$, and
hence Legendre's formula $\eta_1(\tau)-\tau\eta_2(\tau)=2\pi i$ is equivalent to
the transformation formula

\begin{kadisp}{A1.3.12.2}
\frac{\eta_2(-1/\tau)}{\tau} - \eta_2(\tau) = 2\pi i ,
\end{kadisp}

i.e.
\[
\eta_2(-1/\tau) = \tau^2\eta_2(\tau) + 2\pi i\tau
\]
or equivalently $P(-1/\tau)=\tau^2P(\tau)-\frac{6i\tau}{\pi}$.
\end{kathmstar}

\begin{kathmstar}{Remark (A1.3.13).}
Viewing Legendre's relation as saying that
$\langle\omega,\eta\rangle_{\dR}=1$, one can prove it easily using Serre's
cup-product formula, valid on any complete nonsingular curve over $\C$: for any
dfk $\omega$ and any dsk $\eta$, the cup-product
$\langle\eta,\omega\rangle_{\dR}$ is given by the sum
$\Sigma_\cP\operatorname{res}_\cP(f_\cP\cdot\omega)$, where at each point $\cP$,
$f_\cP$ is an element of the $\cP$-adic completion of the function field such
that $\eta=\dd f_\cP$. If one bears in mind that, analytically, we have
$\eta=-\dd\zeta$, then the usual \uline{proof} of Legendre's relation on an
elliptic curve (cf. \cite{whittaker-watson62}, 20.4.11) just becomes an analytic \uline{proof} of
Serre's cup-product formula in that particular case.
\end{kathmstar}

Returning to the relative elliptic curve $\C/\Z+\Z\tau$ over $R$, we have

\begin{kadisp}{A1.3.14}
\begin{pmatrix}\nabla_\tau(\omega) \\ \nabla_\tau(\eta)\end{pmatrix}
= \frac{1}{2\pi i}\begin{pmatrix}
\frac{\pi^2P}{3} & 1 \\
\frac{\pi^4}{9}P^2-\frac{12}{2\pi i}P' & -\frac{\pi^2}{3}P
\end{pmatrix}\begin{pmatrix}\omega \\ \eta\end{pmatrix}
\end{kadisp}

Consider now the differentials $\omega_{\mathrm{can}}=2\pi i\omega$,
$\eta_{\mathrm{can}}=\frac{1}{2\pi i}\eta$, and let
$\theta=\frac{1}{2\pi i}\frac{\dd}{\dd\tau}=q\frac{\dd}{\dd q}$. Then
$\omega_{\mathrm{can}}$ is the canonical differential $\dd t/t$ on the Tate curve\index{Tate curve}
$\mathrm{Tate}(q)$ over $\C((q))$, $\eta_{\mathrm{can}}$ is the d.s.k. "dual" to
$\omega_{\mathrm{can}}$ in the sense of the splitting \eqref{ka:A1.2.6}, and the
Gauss-Manin connection on $H^1_{\dR}(\mathrm{Tate}(q)/\C((q)))$ is given in terms
of $\omega=\frac{1}{2\pi i}\omega_{\mathrm{can}}$ and
$\eta=2\pi i\eta_{\mathrm{can}}$ by

\begin{kadisp}{A1.3.15}
\nabla(\theta)\begin{pmatrix}\omega \\ \eta\end{pmatrix}
&= \frac{1}{2\pi i}\begin{pmatrix}\nabla_\tau(\omega) \\ \nabla_\tau(\eta)\end{pmatrix}
= \frac{-1}{4\pi^2}\begin{pmatrix}
\frac{\pi^2P}{3} & 1 \\
\frac{\pi^4}{9}(P^2-12\theta P) & -\frac{\pi^2P}{3}
\end{pmatrix}\begin{pmatrix}\omega \\ \eta\end{pmatrix} \\
&= \begin{pmatrix}
\frac{-P}{12} & \frac{-1}{4\pi^2} \\
\frac{\pi^2}{36}(P^2-12\theta P) & \frac{P}{12}
\end{pmatrix}\begin{pmatrix}\omega \\ \eta\end{pmatrix}
\end{kadisp}

and hence is given in terms of $\omega_{\mathrm{can}}$, $\eta_{\mathrm{can}}$ by

\begin{kadisp}{A1.3.16}
\nabla(\theta)\begin{pmatrix}\omega_{\mathrm{can}} \\ \eta_{\mathrm{can}}\end{pmatrix}
= \begin{pmatrix}
\frac{-P}{12} & 1 \\
\frac{P^2-12\theta P}{144} & \frac{P}{12}
\end{pmatrix}\begin{pmatrix}\omega_{\mathrm{can}} \\ \eta_{\mathrm{can}}\end{pmatrix} .
\end{kadisp}

\begin{kadisp}{A1.3.17}
\text{The isomorphism } \Omega^1 \simeq \w^{\otimes 2} .
\end{kadisp}

Let $T$ be an arbitrary scheme, $S$ a smooth $T$-scheme, and $E/S$ an elliptic
curve. For any derivation $D\in\Der(S/T)$ and any nowhere vanishing invariant
one-form $\omega$ on $E/S$, we may apply $\nabla(D)$ to $\omega$, view
$\nabla(D)$ as an element of $H^1_{\dR}(E/S)$, and compute the cup-product
$\langle\omega,\nabla(D)\omega\rangle\in\cO_S$. We view this construction as
defining a pairing between $\Der(S/T)$ and $\w^{\otimes 2}$, $\w$ denoting the
line bundle $f_*\Omega^1_{E/S}$ on $S$, or equivalently as a morphism from
$\w^{\otimes 2}$ to $\Omega^1_{S/T}$. The dual mapping
$\Der(S/T)\to(R^1f_*(\widehat{\Omega}^\bullet_E))^{\otimes 2}$ is precisely the
\uline{tangent mapping} of the classifying map from $S$ to the modular stack (or
to the modular scheme\index{modular scheme} $M_n$, if we rigidify the situation with a level $n$
structure). When this map is an isomorphism, the classifying map is
\uline{\'etal\'e}, and we say that $E/S$ is "almost modular".

\begin{kathm}{Corollary}{A1.3.18.}
Consider the Tate curve $\mathrm{Tate}(q)$ over $\Z((q))$. The image of
$\omega_{\mathrm{can}}^{\otimes 2}$ is the differential $\dd q/q$ on $\Z((q))$.
\end{kathm}

\begin{kaproof}
The assertion is that
$\langle\omega_{\mathrm{can}},\nabla(\theta)\omega_{\mathrm{can}}\rangle=1$. It
suffices to check this over $\C((q))$, where we have
$\nabla(\theta)(\omega_{\mathrm{can}})=\frac{-P}{12}\omega_{\mathrm{can}}+\eta_{\mathrm{can}}$.
As $\langle\omega_{\mathrm{can}},\omega_{\mathrm{can}}\rangle=0$, and
$\langle\omega_{\mathrm{can}},\eta_{\mathrm{can}}\rangle=1$.
\end{kaproof}

\sechead{A1.4}{The Gauss-Manin connection and Serre's $\partial$ operator (\cite{serre73}): d'apr\`es Deligne}

A series $f(q)\in\C[[q]]$ is (the q-expansion of) a modular form\index{modular form} of weight $k$
if and only if $f(q)\cdot(\omega_{\mathrm{can}})^{\otimes k}$ extends to a
"global" section of $\w^{\otimes k}$, i.e. one which is "defined" for all
families of elliptic curves\index{elliptic curve} $/\C$, or equivalently if there exist integers $a,b$
with $a-b=k$ such that
$f(q)\cdot(\omega_{\mathrm{can}})^{\otimes a}\cdot(\eta_{\mathrm{can}})^{\otimes b}$
extends to a "global" section of $\Sym^{a+b}(H^1_{\dR})$, in the same sense.

We now view the Gauss-Manin connection\index{Gauss--Manin connection} on $H^1_{\dR}(E/S)$, where $S$ is a
smooth $T$-scheme, as an arrow
$\nabla\colon H^1_{\dR}(E/S)\longrightarrow H^1_{\dR}(E/S)\tens\Omega^1_{S/T}$.
Its $k$'th symmetric power is a connection on $\Sym^k(H^1)$, so an arrow

\[
\Sym^k(H^1) \longrightarrow \Sym^k(H^1)\tens\Omega^1_{S/T} .
\]

If $E/S$ is "almost" modular, we have isomorphism
$\Omega^1_{S/T}\sim\w^{\otimes 2}$, so we may view this last arrow as an arrow

\[
\Sym^k(H^1) \longrightarrow \Sym^k(H^1)\tens\w^{\otimes 2} .
\]

Suppose now that $6=2\cdot 3$ is \uline{invertible} in $S$. Then we have a
splitting $H^1_{\dR}(E/S)\sim\w\oplus\w^{-1}$, whose $k$'th symmetric power is a
splitting $\Sym^k(H^1_{\dR}(E/S))\sim{}\allowbreak\sum_{j=0}^{k}\w^{\otimes k-2j}$, and we may
interpret the Gauss-Manin connection as an arrow

\begin{kadispb}{A1.4.1}
\sum_{j=0}^{k}\w^{\otimes k-2j} \longrightarrow
\sum_{j=0}^{k}\w^{\otimes k-2j}\tens\w^{\otimes 2}
= \sum_{j=0}^{k}\w^{\otimes k+2-2j} .
\end{kadispb}

Suppose that $f$ is the q-expansion of a modular form of weight $k$. Then for
any integers $a$ and $b$ such that $a-b=k$,
$f\cdot(\omega_{\mathrm{can}})^{\otimes a}\tens(\eta_{\mathrm{can}})^{\otimes b}$
extends to a global section of $\Sym^{a+b}(H^1)$. Hence its \uline{image} under
the Gauss-Manin connection extends to a global section of
$\Sym^{a+b}(H^1)\tens\w^{\otimes 2}$. But its image under Gauss-Manin is

\[
\begin{aligned}
&\theta(f)\cdot(\omega_{\mathrm{can}})^{\otimes 2}\cdot(\omega_{\mathrm{can}})^{\otimes a}\cdot(\eta_{\mathrm{can}})^{\otimes b} \\
&+ f\cdot a\cdot(\omega_{\mathrm{can}})^{\otimes a-1}\left(\frac{-P}{12}\omega_{\mathrm{can}}+\eta_{\mathrm{can}}\right)
\tens(\omega_{\mathrm{can}})^{\otimes 2}\tens(\eta_{\mathrm{can}})^b \\
&+ f\cdot(\omega_{\mathrm{can}})^{\otimes a}\cdot b\cdot(\eta_{\mathrm{can}})^{\otimes b-1}
\left(\frac{P^2-12\theta P}{144}\omega_{\mathrm{can}}+\frac{P}{12}\eta_{\mathrm{can}}\right)\cdot(\omega_{\mathrm{can}})^{\otimes 2}
\end{aligned}
\]

which we group according to the decomposition
$\Sym^{a+b}(H^1)\tens\w^{\otimes 2}\simeq\sum_{j=0}^{a+b}\w^{a+b+2-2j}$,

\[
\begin{aligned}
= &\left\{\theta(f)-(a-b)\cdot f\cdot\frac{P}{12}\right\}\cdot(\omega_{\mathrm{can}})^{\otimes a+2}\cdot(\eta_{\mathrm{can}})^{\otimes b} \\
&+ \{af\}\cdot(\omega_{\mathrm{can}})^{\otimes a+1}\cdot(\eta_{\mathrm{can}})^{\otimes b+1} \\
&+ \left\{bf\cdot\frac{P^2-12\theta P}{144}\right\}(\omega_{\mathrm{can}})^{\otimes a+3}\cdot(\eta_{\mathrm{can}})^{\otimes b-1} .
\end{aligned}
\]

Thus we conclude that if $f$ is modular of weight $k=a-b$, then

\begin{kadisp}{A1.4.2}
\begin{cases}
\theta(f)-kf\cdot\frac{P}{12} & \text{is modular of weight } k+2 = a+2-b \\
af & \text{is modular of weight } k = a-b \\
b\cdot f\cdot\left[\frac{P^2-12\theta P}{144}\right] & \text{is modular of weight } k+4 = a+3-(b-1) .
\end{cases}
\end{kadisp}

[Serre's $\partial$ operator is $\partial(f)=12\,\theta(f)-k\cdot P\cdot f$ for
$f$ modular of weight $k$, hence $\partial f$ is modular of weight $k+2$.]

\begin{kathm}{Corollary}{A1.4.3.}
$P^2-12\,\theta P$ is modular of weight 4, hence
$\underline{P^2-12\,\theta P = E_4 \xlongequal{\text{dfn}} Q}$.
\end{kathm}

\begin{kaproof}
Take $f=1$ which is modular of weight $0=1-1$, to see that $P^2-12\,\theta P$ is
modular of weight 4. As it has constant term 1, it is necessarily $E_4$.
\end{kaproof}

\begin{kathm}{Corollary}{A1.4.4.}
(Deligne) $P=\frac{\theta\Delta}{\Delta}$, where $\Delta$ denotes the unique
normalized cusp form of weight 12, the discriminant $(E_4^3-E_6^2)/1728$.
\end{kathm}

\begin{kaproof}
$\theta(\Delta)-\Delta\cdot P$ is a \uline{cusp form} of weight 14 and level 1,
and there are none save zero.
\end{kaproof}

\begin{kathm}{Corollary}{A1.4.5.}
The Gauss-Manin connection on $H^1_{\dR}$ of $\mathrm{Tate}(q)$ over
$\Z[1/6]((q))$ is given by
\end{kathm}

\begin{kadisp}{A1.4.6}
\begin{pmatrix}\nabla(\theta)(\omega_{\mathrm{can}}) \\ \nabla(\theta)(\eta_{\mathrm{can}})\end{pmatrix}
= \begin{pmatrix}\frac{-P}{12} & 1 \\ \frac{Q}{144} & \frac{P}{12}\end{pmatrix}
\begin{pmatrix}\omega_{\mathrm{can}} \\ \eta_{\mathrm{can}}\end{pmatrix} .
\end{kadisp}

\begin{kaproof}
$\omega_{\mathrm{can}}$, $\eta_{\mathrm{can}}$ give a base of $H^1_{\dR}$ over
$\Z(\frac{1}{6})((q))\subset\C((q))$, and we have the desired assertion by
transcendental means over $\C((q))$.
\end{kaproof}

\begin{kathmstar}{Remarks.}
\end{kathmstar}

\kaitemb[1-c]{1.} The \uline{value at 0} of the connection matrix is
$\begin{pmatrix}-1/12 & 1 \\ 1/144 & 1/12\end{pmatrix}$, which is a
\uline{nilpotent} matrix. This shows that the \uline{canonical extension} (in the
sense of (\cite{deligne70})) of $H^1_{\dR}$ with its Gauss-Manin connection to $\infty$ is
given by the free module with base $\omega_{\mathrm{can}}$, $\eta_{\mathrm{can}}$.

\kaitemb[2-c]{2.} We have $(\nabla(\theta))^2(\omega_{\mathrm{can}})=0$ (because the
\uline{periods} of $\omega$ are 1 and $\tau$ both killed by
$\left(\frac{\dd}{\dd\tau}\right)^2$), hence by Igusa \cite{igusa58}, the
\uline{Hasse-invariant} has a q-expansion $f(q)\in\F_p[[q]]$ which satisfies
$\theta^2f=0$, so writing $f=\Sigma\,a_nq^n$, we have $(a_n)^2=0$, hence $a_n=0$,
hence $f=a_0+a_pq^p+\dots$. By \uline{direct calculation} (of the coefficient of
$x^{p-1}$ in the $\frac{p-1}{2}$'th power of
$4x^3-\frac{E_4(0)}{12}x+\frac{E_6(0)}{216}=4x^3-\frac{x}{12}+\frac{1}{216}$,
(cf. \cite{katz72}, 2.3.7.14)), we compute $a_0=1$, hence $f\equiv 1\bmod(q^p)$. As the
same is also true for the reduction mod p of $E_{p-1}$, we have
$E_{p-1}\equiv f\bmod(p,q^p)$, hence $E_{p-1}-f$ is a cusp form mod p of weight
$p-1$ and level 1 with a zero of order $\geq p$, hence \uline{vanishes} mod p.
Thus $E_{p-1}$ mod p is the Hasse invariant\index{Hasse invariant}, and $f(q)$ is identically 1. (We
gave Deligne's original and more conceptual proof of this fact in \eqref{ka:2.1}.)

\sechead{(A1.5)}{Formulas}

\kaitem{A1.5.0} For $n\geq 3$, $\Mbar_n$ is proper and smooth over
$\Z[1/n]$, and its inverse image over $\Z[1/n,\zeta_n]$ is the disjoint union of
$\varphi(n)$ proper smooth schemes with geometrically connected fibres
$\Mbar_n^\zeta$, one for each primitive $n$'th root of unity $\zeta$
(corresponding to the value of the e.m. pairing on the given basis of points of
order $n$). The $\Z[1/n,\zeta_n]$ schemes $\Mbar_n^\zeta$ are non-canonically
isomorphic to each other. We give below the formulas for their (common) genus,
the (common) number of their cusps, and the \uline{degree} of the invertible
sheaf $\w$. The method of deducing such relations is very simple: one notes that
by flatness, the fact that $\Mbar_n^\zeta-\Mbar_n^s$ is a disjoint union of
sections, and the isomorphism
$\w^{\otimes 2}\simeq\Omega^1_{\Mbar_n^\zeta/\Z[1/n,\zeta_n]}$ (log "cusps"), it
suffices to calculate these invariants for any geometric fibre
$\Mbar_n^\zeta\tens k$ [$k$ \uline{any} algebraically closed field containing
$1/n$]. One then applies the standard Hurwitz formula to the morphism
$\Mbar_n^\zeta\tens k\longrightarrow\PP^1_k$ provided by the j-invariant. A
closed point of $\PP^1_k$ other than $\infty$ "is" an elliptic curve\index{elliptic curve} $E$ over
$k$, up to isomorphism. The points of $\Mbar_n^\zeta\tens k$ lying over it are
the set of all level $n$ structures on $E$ such that the value of the e.m.
pairing on the given basis of ${}_nE$ is $\zeta$, \uline{modulo} the natural
action of $\Aut(E)$ of ${}_nE$. The \uline{cardinality} of the fibre over the
point "E" is thus $\#SL_2(\Z/n\Z)/\#\Aut(E)$. For $j(E)\neq 0$, 1728,
$\Aut(E)=\pm 1$, and hence over $\PP^1_k-\{0,1728,\infty\}$, the projection is
\'etal\'e of degree $\#SL_2(\Z/n\Z)/2$. The fibre over 0 has
$\#SL_2(\Z/n\Z)/6$ points, and that over 1728 has $\#SL_2(\Z/n\Z)/4$ points. The
points over $\infty$ are the cusps, each of which is ramified of degree $n$,
hence the number of cusps is $\#SL_2(\Z/n\Z)/2n$. Letting $\chi$ denote the
\uline{topological} Euler characteristic, we thus have the formula:

\[
\chi(\Mbar_n^\zeta\tens k) = \#SL_2(\Z/n\Z)\left[\frac{1}{6}+\frac{1}{4}+\frac{1}{2n}\right]
+ \#SL_2(\Z/n\Z)[1/2]\cdot\chi(\PP^1-\{0,1728,\infty\}) ,
\]

i.e.,
$\chi(\Mbar_n^\zeta\tens k)=\#SL_2(\Z/n\Z)\left[\frac{1}{6}+\frac{1}{4}+\frac{1}{2n}-\frac{1}{2}\right]
=\#SL_2(\Z/n\Z)\cdot\left[\frac{6-n}{12n}\right]$. Now

\[
\#SL_2(\Z/n\Z) = n^3\prod_{p\mid n}\left(1-\frac{1}{p^2}\right) ,
\]

so we have finally

\kaitem{A1.5.1} \uline{Formulas}

\begin{kadisp}{A1.5.2}
1-\text{genus}(\Mbar_n^\zeta) = \frac{6-n}{24n}\#SL_2(\Z/n\Z)
= \frac{n^2(6-n)}{24}\prod_{p\mid n}\left(1-\frac{1}{p^2}\right) ;
\end{kadisp}

\begin{kadisp}{A1.5.3}
\#\text{cusps on }\Mbar_n^\zeta = \frac{1}{2n}\#SL_2(\Z/n\Z)
= \frac{n^2}{2}\prod_{p\mid n}\left(1-\frac{1}{p^2}\right) ;
\end{kadisp}

\begin{kadisp}{A1.5.4}
\text{degree}(\w)\text{ on }\Mbar_n^\zeta
&= \frac{1}{2}\text{deg}(\Omega^1(\log\text{ cusps})) = \frac{1}{2}(2g-2+\#\text{cusps}) \\
&= \left(\frac{n-6}{24n}+\frac{1}{4n}\right)\#SL_2(\Z/n\Z) \\
&= \frac{1}{24}\#SL_2(\Z/n\Z) .
\end{kadisp}

\kaitem{A1.5.5} \uline{Sample consequences}

$\Mbar_n^\zeta$ has genus zero only for $n=3,4,5$, and genus one only for $n=6$.
We always have $\deg(\w^{\otimes 2})>2g-2$, but $\deg(\w)>2g-2$ \uline{only} for
$3\leq n\leq 11$. For $n=3,4,5$, $\Mbar_n^\zeta$ is a $\PP^1$, hence $\w$ is
uniquely determined by its degree; $\w=\cO(1)$ on $\Mbar_3^\zeta$,
$\w=\cO(2)$ on $\Mbar_4^\zeta$, $\w=\cO(5)$ on $\Mbar_5^\zeta$.
